\frenchspacing
\documentclass{amsart}
\usepackage{amssymb}
\usepackage{fancyhdr}
\usepackage{bbm}
\usepackage{enumerate}
\usepackage{graphicx}
\usepackage{tikz}
\usepackage{minted}
\pagestyle{fancy}
\usepackage[margin=1in]{geometry}
\newgeometry{left=1.5cm, right=1.5cm, top = 1.5cm}
\fancyhfoffset[E,O]{0pt}
\allowdisplaybreaks


\rhead{Andrew Lys}   %% <-- your name here
\chead{Problem}
\cfoot{\thepage}
\lhead{\today}



%% your macros -->
\newcommand{\nn}{\mathbb N}    %% naturals
\newcommand{\zz}{\mathbb Z}    %%integers
\newcommand{\rr}{\mathbb R}    %% real numbers
\newcommand{\cc}{\mathbb C}    %% complex numbers
\newcommand{\ff}{\mathbb F}
\newcommand{\qq}{\mathbb Q}
\newcommand{\cA}{\mathcal{A}}
\newcommand{\cP}{\mathcal{P}}
\newcommand{\cN}{\mathcal{N}}
\newcommand{\cM}{\mathcal{M}}
\newcommand{\cB}{\mathcal{B}}
\newcommand{\cL}{\mathcal{L}}
\newcommand{\limn}{\lim_{n \to \infty}} %%lim n to infty shorthand
\newcommand{\va}{\mathbf{a}}
\newcommand{\vb}{\mathbf{b}}
\newcommand{\vc}{\mathbf{c}}
\newcommand{\vx}{\mathbf{x}}
\newcommand{\vy}{\mathbf{y}}
\newcommand{\vv}{\mathbf{v}}
\newcommand{\vw}{\mathbf{w}}
\newcommand{\vu}{\mathbf{u}}
% Source - https://tex.stackexchange.com/q/79434
% Posted by Memming, modified by community. See post 'Timeline' for change history
% Retrieved 2026-02-11, License - CC BY-SA 3.0

\newcommand\ind{\protect\mathpalette{\protect\independenT}{\perp}}
\def\independenT#1#2{\mathrel{\rlap{$#1#2$}\mkern2mu{#1#2}}}

\DeclareMathOperator{\Var}{\mathbb{V}ar}  %% variance
\DeclareMathOperator{\Aut}{Aut}
\DeclareMathOperator{\Sym}{Sym}
\DeclareMathOperator{\Cov}{\mathbb{Cov}}
\DeclareMathOperator{\ee}{\mathbb{E}}
\newtheorem{theorem}{Theorem}
\theoremstyle{definition}
\newtheorem{definition}{Definition}
\newtheorem{exercise}{Exercise}


\begin{document}
\noindent
Problem Set 4  \hfill \today  %% <-- Update Notes here ***
\smallskip
\hrule
\smallskip
\noindent
Solutions by {\bf Andrew Lys} \qquad   %% <-- your name here ***
  {\tt alys(at)ncsu.edu}

\vspace{0.5cm}
\section*{Problem 1. Transformers as Dependency Parsers}
\subsection*{Part (a)} 
\begin{proof}
    \[
    \alpha(\mathbf{dobj}, \mathbf{prefer}, w) = \delta_{w, \text{flight}}
    \]
\end{proof}
\subsection*{Part (b)}
\begin{proof}
    We are interested in penalizing attention vectors that spread attention across multiple words, 
    and rewarding attention vectors that focus on a single word. 
    A natural candidate for this is the shannon entropy of the attention vector. Let $p$ be a probability vector. Then
    \[
    H(p) := -\sum_{i} p_i \log p_i
    \] 
    Since we wish to penalize all heads, we have the following $\cL_{\text{focus}}$ definition:
    \[
    \cL_{\text{focus}} := \sum_{\ell, k, t_1} H(\alpha[\ell, k, t_1, T_2])
    \]
    Additionally, we add a hyperparameter $\lambda$ to control the strength of the regularization.
    \[
    \cL = \cL_{\text{LM}} + \lambda \cL_{\text{focus}}
    \]
\end{proof}

\subsection*{Part (c)}
\begin{proof}
    A natural candidate for penalizing distance from $t_1$ is just the absolute distance from $t_1$, scaled by the attention weight. 
    Squared distance is probably a better choice. 
    Additionally, in order to cap the maximum penalty, to allow for the possibility of outliers, we introduce a hyperparameter $L_{\text{max}}$.
    As before, we introduce a hyperparameter $\lambda$ to control the strength of the regularization.
    \[
    \cL_{\text{near}} := \sum_{\ell, k, t_1} \ee^{\alpha[\ell, k, t_1, T_2]}[(t_1 - t_2)^2 \land L_{\text{max}}]
    \]
\end{proof}

\subsection*{Part (d)}
The universality assumption is that the model will learn the optimal parameters for the language modeling task,
so loss shaping will change the optimal parameters. 
In the real world, it is unlikely that the model will learn the optimal parameters, so loss shaping helps guide the model towards better parameters.

\section*{Problem 2. Adjusting Temperature for Dimension}
Consider adding a temperature parameter to the softmax:
\[
\alpha[\ell, t, w] = \underset{w}{\mathrm{softmax}} \beta h[\ell, t, I] e[w, I]
\]
\subsection*{Part (a)} 
\begin{proof}
    \begin{align*}
        \ee[\beta h[t, I] e[w, I]] &= 0 \\
        \Var(\beta h[t, I] e[w, I]) &= \beta^2 \Var(h[t, I] e[w, I])  = \beta^2 \sum_i \ee[h[t, i]^2] \ee[e[w, i]^2] = \beta^2 I
    \end{align*}
    Therefore we let $\beta = I^{-1/2}$. 
\end{proof}
\subsection*{Part (b)}
\begin{proof}
    This is exactly the normalization used in the softmax. 
\end{proof}

\section*{Parameterizing Inner-Product}

\subsection*{Part (a)}
We require that we have multiple heads, so we need $k$ to be a parameter. 
\subsection*{Part (b)}
\begin{align*}
    P(y | T_2) &= \underset{t_2}{\mathrm{softmax}}\,\frac{1}{\sqrt{I}} L[\ell, t_1, J]^\top W^Q_{\ell + 1}[k, J, I] W^K_{\ell + 1}[k, I, J] L[\ell, t_2, J] 
\end{align*}

\section*{Problem 4. Repetion at Low Temperatures}
Let $\beta = \infty$. Then we have:
\[
\Pr(w_{t + 2} | w_{t}, w_{t + 1}) = \lim_{\beta \to \infty} \underset{w_{t + 2}}{\mathrm{softmax}}\,\beta h[t + 2, I] e[w_{t + 2}, I]
\]
which as seen in the first homework, is a one-hot distribution on the word with the highest inner product with $h[t + 2, I]$.
Each word is deterministic, so whenever we encounter a repeated pair of words, we will always repeat the same word, leading to infinite repetition.

\section*{Problem 5. Eliminating the Key Matrix}
Self attention is computed by the following equations:
\begin{align*}
    \mathrm{Query}_{\ell + 1} [k, t, i] &= W^Q_{\ell + 1}[k, i, J] L_{\ell}[t, J] \\
    \mathrm{Key}_{\ell + 1} [k, t, i] &= W^K_{\ell + 1}[k, i, J] L_{\ell}[t, J] \\
    \alpha_{\ell + 1} [k, t_1, t_2] &= \underset{t_2}{\mathrm{softmax}}\, \left[
        \frac1{\sqrt{I}} \mathrm{Query}_{\ell + 1} [k, t_1, I] \mathrm{Key}_{\ell + 1} [k, t_2, I]
    \right]
\end{align*}
$W^Q$ and $W^K$ are of the same shape $[K, I, J]$. Typically $I < J$
\end{document}

